% Compile-ready section. Requires: booktabs, tabularx, array, enumitem.

\section{GlyphChain Assurance and Privacy Boundary}
\label{sec:glyphchain-assurance-privacy-boundary}

This section defines the assurance language that Tessaris may display for a
Legal document proof, the information permitted on GlyphChain and the claims
that remain prohibited.  These terms are normative product and release
boundaries.  User-interface wording, sales material and exported evidence must
not collapse a weaker state into a stronger one.

\subsection{Supported Assurance Levels}

\begin{tabularx}{\textwidth}{@{}
  >{\raggedright\arraybackslash\bfseries}p{0.22\textwidth}
  >{\raggedright\arraybackslash}X
  >{\raggedright\arraybackslash}X@{}}
\toprule
Displayed state & What Tessaris has verified & What it does not claim \\
\midrule
Local evidence created &
The document, signer commitments and audit history have been hashed and signed
by the originating Tessaris device. &
Network publication, delivery, legal enforceability or quorum finality. \\
\addlinespace

Ledger committed --- quorum pending &
The exact proof transaction is stored in a GlyphChain ledger block and its
Merkle inclusion path verifies. &
Agreement by a validator quorum. \\
\addlinespace

Signatures valid --- trust anchor not supplied &
The certificate's validator membership, Ed25519 votes and quorum power are
internally valid. &
That the validator set is the legitimate GlyphChain set. \\
\addlinespace

Trust-anchored quorum final &
The version~2 certificate is valid against a separately supplied checkpoint and
signs the exact ledger commitment and transaction root containing the proof. &
Legal validity, qualified electronic-signature status, identity assurance
beyond the configured signing policy, or uninterrupted public-network
availability. \\
\bottomrule
\end{tabularx}

The application must not shorten these distinct states to a generic
``blockchain certified'' label.  A ledger commitment is not quorum finality,
and a certificate carried inside a receipt cannot appoint its own trust anchor.

\subsection{Data Placed on GlyphChain}

GlyphChain stores commitments and protocol metadata.  The document body,
signer name, signer email address, message body, delivery attachment and private
Legal matter record are not written into the proof transaction.  The
transaction binds opaque identifiers and hashes for the document, signed-copy
manifest, ordered signer commitments and terminal audit event.

Anyone who obtains the private document and portable receipt can test whether
the bytes match the public commitment.  A commitment is not encryption and
must not be presented as making a weak or predictable document secret.

\subsection{Local Data and Key Separation}

\begin{itemize}[leftmargin=*]
  \item Legal matter and envelope records use AES--256--GCM encryption and
        production Keychain custody.
  \item The Legal proof device identity is a dedicated Ed25519 key and is not
        the validator-finality identity.
  \item Peer-to-peer transport authentication uses a separate Ed25519 key.
  \item Validator quorum-finality votes use a separate Ed25519 key anchored in
        the validator checkpoint.  Startup rejects reuse of the peer-to-peer
        transport key as the finality key.
  \item A normal desktop is a light node.  Relay and validator duties require
        explicit local approval; installing Tessaris does not silently
        contribute public bandwidth or consensus votes.
\end{itemize}

\subsection{Retention and Deletion}

Local Legal records follow the configured retention schedule and applicable
Legal holds.  Deletion requires a portable export and an exact confirmation,
then leaves an encrypted, content-free destruction receipt.  Deleting local
content does not erase a public commitment or transaction hash from a
distributed ledger.  Tessaris must disclose that distinction before proof
submission.

\subsection{Claims That Remain Prohibited}

Until separately assessed and approved, Tessaris must not claim:

\begin{itemize}[leftmargin=*]
  \item qualified electronic-signature or qualified trust-service status;
  \item guaranteed Legal enforceability in every jurisdiction;
  \item anonymous or untraceable network use;
  \item permanent availability solely because a proof was submitted;
  \item production-grade decentralisation before the multi-device network
        gates pass;
  \item that a local development checkpoint is an independently distributed
        trust anchor; or
  \item that Photon credits or currency are required for Legal document proof.
\end{itemize}

\subsection{Production Publication Gate}

The following artefacts must be versioned with every production release:

\begin{enumerate}[leftmargin=*]
  \item this assurance and privacy-boundary section;
  \item the network-participation guide;
  \item the supported retention and recovery commitments;
  \item the current public validator-checkpoint file; and
  \item the checkpoint identifier or digest published through a genuinely
        independent channel.
\end{enumerate}

The validator checkpoint and its trust pin must not be obtained solely from the
same mutable source.  Production publication remains gated until a clean-device
verification succeeds using the versioned checkpoint, independently published
pin, private evidence package and standalone verifier.
